%% 26_emg.tex -- Chapter 26: EMG and Neural Spike Processing
%% Status: OUTLINE (build stage 5). Replace the outline with the full chapter when it is written.
\chapter{EMG and Neural Spike Processing}\label{ch:emg}
\chapterlead{Muscle signals are read through their envelope; neural signals through the timing of individual spikes.}

\begin{specnote}[Chapter status: outline]
This chapter is an outline. It will be written in build stage 5 (Part IV, including bioimpedance, with Labs 11--13 and 15). The topics, the design question and the labs below are fixed; the derivations, worked examples and figures are still to come.
\end{specnote}

\begin{objectives}
\begin{itemize}
  \item compute EMG envelopes and the time-domain features MAV, WL, ZC and SSC;
  \item classify hand gestures on NinaPro;
  \item detect neural spikes with the nonlinear energy operator (NEO);
  \item explain spike sorting and why implants compress before transmitting.
\end{itemize}
\end{objectives}

\begin{insight}[Design question]
How much resolution, sample rate and channel count does gesture recognition or spike detection actually need?
\end{insight}

\section{EMG envelopes and features}\label{sec:ch26-1}
Rectification and smoothing; mean absolute value, waveform length, zero crossings and slope-sign changes.

\section{Gesture classification}\label{sec:ch26-2}
LDA and SVM classifiers on the NinaPro database~\cite{ninapro}.

\section{Spike detection}\label{sec:ch26-3}
Thresholds and the NEO; detection versus the noise of \cref{ch:noise}.

\section{Spike sorting and compression}\label{sec:ch26-4}
Clustering waveforms; why a thousand-channel implant must send spikes, not samples.

\begin{tryit}[Labs for this chapter]
Lab~13 (\cref{sec:lab13}) computes MAV, WL, ZC and SSC on NinaPro, classifies with LDA or SVM and repeats with fixed-point features.
\end{tryit}

\paragraph{Read with this chapter} Atzori et al.~\cite{ninapro}.
